    %%%%%%%%%%%%%%%%%%%%%%%%%%%%%%%%%%%%%%%%%%%%%%%%%%%%%%%%%%%%%%%%%%%%%%%%%%%%%%%%%%%%%%%%%%%%%%%%%%%%%%%%%%%%%%
    %%%%%%%%%%%%%%%%%%%%%%%%%%%%%%%%%%%%%%%%%%%%%%%%% PCS I10 %%%%%%%%%%%%%%%%%%%%%%%%%%%%%%%%%%%%%%%%%%%%%%%%%%
    %%%%%%%%%%%%%%%%%%%%%%%%%%%%%%%%%%%%%%%%%%%%%%%%%%%%%%%%%%%%%%%%%%%%%%%%%%%%%%%%%%%%%%%%%%%%%%%%%%%%%%%%%%%%%%
    {{ waterMarkText }}

    \renewcommand{\DTRPcs}{Islanding} % DTR pcs definition
    \renewcommand{\PCSname}{PCS_RTE-I16z3}
    \renewcommand{\OCname}{DeltaP10DeltaQ4Injection}


    \subsection{Test I10 - Islanding event in injection}

    Checks for compliant behavior of the generator under an
    \emph{islanding} event grid scenario, i.e., the network has been split in
    such a way that the generator is left as the only synchronous machine
    that should sustain the frequency of its subnetwork. The generator
    should maintain the stability until the restoration process
    re-synchronizes this subnetwork to the bulk transmission network.

    \GridCircuitZthree

    The separated network is modeled as the DTR describes it. A load of type
    \code{LoadAlphaBeta} with alpha = beta = 0, so that it depends neither on
    voltage nor on frequency, is attached to the PDR bus, and the PDR bus is
    connected through a line of reactance $b$ to an inertial grid that stands
    for the fictitious synchronous condenser: it holds a constant voltage, has
    an inertia constant of 1 s on $P_{max\_unite}$, no damping and practically
    no governor response, and carries no power before the event. The load
    starts at the initial operating point of the plant and, at the event, its
    consumption steps by 10\% of $P_{max\_unite}$ in active power and by 4\%
    of $P_{max\_unite}$ in reactive power.

    \begin{description}
        \item Initial conditions used at the PDR bus:
        \begin{itemize}
            \item $P = 0.8 P_{max\_injection\_unite}$
            \item $Q = 0$
            \item $U = U_{dim}$
            \item $f = 50 Hz$ (this is always the default in Dynawo)
        \end{itemize}
    \end{description}

    \subsubsection{Simulation parameters}

    Solver and parameters used in the simulation:
    \begin{center}
        \begin{tabular}{lc}
            \toprule
           \textbf{Parameter} & \textbf{Value (default)} \\
            \midrule
            \BLOCK{for row in solverPCSI16z3IslandingDeltaP10DeltaQ4Injection}
            {{row[0]}}         & {{row[1]}}                         \\
            \BLOCK{endfor}
            \bottomrule
        \end{tabular}
    \end{center}

    \GridSimulationHeaderZthree{
        \GridCurvesDescriptionZthree
        \GridCurvesIslandingZthreeDescriptionZthree
        \GridCurvesFrequencyZthreeDescriptionZthree
    }
    \GridCurvesZthree
    \GridCurvesIslandingZthree
    \GridCurvesFrequencyZthree
    \\[2\baselineskip]
    Go to  {{ linkPCSI16z3IslandingDeltaP10DeltaQ4Injection }}

    \subsubsection{Analysis of results}

    \noindent Dynamic behaviour: {{ stabilizedPCSI16z3IslandingDeltaP10DeltaQ4Injection }}

    \noindent Analysis of results for PCS \DTRPcs. Values extracted
    from the simulation:
%\noindent No analysis is needed in PCS \DTRPcs.
    \begin{center}
        \scriptsize
        \begin{tabular}{@{}lcccccc@{}}
            \toprule
            & \multicolumn{3}{c}{Pre-event} & \multicolumn{3}{c}{Event} \\
            \cmidrule(lr){2-4}\cmidrule(lr){5-7}
            & {MXE}      & {ME}       & {MAE}      & {MXE}      & {ME}       & {MAE}      \\
            \midrule
            \BLOCK{for row in rmPCSI16z3IslandingDeltaP10DeltaQ4Injection}
            {{row[0]}} & {{row[1]}} & {{row[2]}} & {{row[3]}} & {{row[4]}} & {{row[5]}} & {{row[6]}} \\
            \BLOCK{endfor}
            \bottomrule
        \end{tabular}
    \end{center}

    \subsubsection{Compliance checks}

    {{ missedColumns }}

    \noindent Compliance thresholds on the curves:
    \begin{center}
        \scriptsize
        \begin{tabular}{@{}lcccccc@{}}
            \toprule
            & \multicolumn{3}{c}{Pre-event} & \multicolumn{3}{c}{Event} \\
            \cmidrule(lr){2-4}\cmidrule(lr){5-7}
            & {MXE}      & {ME}       & {MAE}      & {MXE}      & {ME}       & {MAE}      \\
            \midrule
            \BLOCK{for row in thmPCSI16z3IslandingDeltaP10DeltaQ4Injection}
            {{row[0]}} & {{row[1]}} & {{row[2]}} & {{row[3]}} & {{row[4]}} & {{row[5]}} & {{row[6]}} \\
            \BLOCK{endfor}
            \bottomrule
        \end{tabular}
    \end{center}

    \noindent Compliance checks on the curves:
    \begin{center}
        \scriptsize
        \begin{tabular}{@{}lccccccc@{}}
            \toprule
            & \multicolumn{3}{c}{Pre-event} & \multicolumn{3}{c}{Event} & \\
            \cmidrule(lr){2-4}\cmidrule(lr){5-7}
            & {MXE}      & {ME}       & {MAE}      & {MXE}      & {ME}       & {MAE}      & Compl.     \\
            \midrule
            \BLOCK{for row in emPCSI16z3IslandingDeltaP10DeltaQ4Injection}
            {{row[0]}} & {{row[1]}} & {{row[2]}} & {{row[3]}} & {{row[4]}} & {{row[5]}} & {{row[6]}} & {{row[7]}} \\
            \BLOCK{endfor}
            \bottomrule
        \end{tabular}
    \end{center}

    \noindent In steady state after the event, the absolute average error must not exceed {{steadystatethreshold}}\% (configured value):
    \begin{center}
        \scriptsize
        \begin{tabular}{cllc}
            \toprule
            Variable   & MAE        & Final Error & Compliance \\
            \midrule
            \BLOCK{for row in ssemPCSI16z3IslandingDeltaP10DeltaQ4Injection}
            {{row[0]}} & {{row[1]}} & {{row[2]}}  & {{row[3]}} \\
            \BLOCK{endfor}
            \bottomrule
        \end{tabular}
    \end{center}
